\documentclass[conference]{IEEEtran}

% Paquetes necesarios
\usepackage[utf8]{inputenc}
\usepackage[spanish]{babel}
\usepackage{amsmath}
\usepackage{amsfonts}
\usepackage{amssymb}
\usepackage{graphicx}
\usepackage{cite}
\usepackage{xcolor}
\usepackage{array}
\usepackage{multirow}
\usepackage{multicol}
\usepackage{url}
\usepackage{fancyhdr}
\usepackage{watermark}
\usepackage{subfig}

%\documentclass{article} % o la clase que uses

% Paquetes para bibliografía
\usepackage[utf8]{inputenc}
\usepackage[spanish]{babel} % si escribes en español
\usepackage[hyphens]{url}
\usepackage[breaklinks=true]{hyperref}
\usepackage[backend=biber,style=numeric,sorting=none]{biblatex}

% Especificar el archivo de bibliografía
\addbibresource{referencias.bib}


% Configuración de idioma y codificación
\selectlanguage{spanish}

% Información del documento
\title{Introducción a la Investigación:\\Tarea Corta 1}

\author{
\IEEEauthorblockN{Emmanuel Barrantes Vargas}\
\IEEEauthorblockN{Isaac Francisco Moreno Fuentes}\\
\IEEEauthorblockA{Tecnológico de Costa Rica\\
Maestría en Ciencia de la Computación}
}

\begin{document}

% Título
\maketitle

% Abstract
\begin{abstract}
Este documento presenta el resultado de la investigación realizada alrededor de los conceptos de Revisión Sistemática de Literatura (Systematic Literature Review) y Estudio de Mapeo Sistemático (Systematic Mapping Study), su definición, principales diferencias, y explicación básica de ambos procesos.
\end{abstract}

% Palabras clave
\begin{IEEEkeywords}
Revisión Sistemática de Literatura, Revisión Sistemática de Literatura
\end{IEEEkeywords}

% Introducción
%\section{Definiciones}

% Incluir todas las secciones
\section{Estudio de Mapeo Sistemático}

Un Estudio de Mapeo Sistemático (o SMS por sus siglas en inglés Systematic Mapping System) es un tipo de estudio cuyo objetivo es proporcionar una visión general de un campo de investigación para identificar qué evidencias están disponibles sobre un tema. \cite{definicion_sms} Se enfoca en la amplitud y no en la profundidad del área investigada. Ademas, permite identificar tendencias de investigación a lo largo del tiempo, y descubrir posibles brechas y puntos de enfoque, proporcionando una síntesis de la estructura del tema elegido mediante la clasificación y categorización de los estudios.
\section{Revisión Sistemática de Literatura}
Una Revisión Sistemática de Literatura (o SLR por sus siglas en inglés Systematic Literature Review) es un medio para identificar, evaluar e interpretar toda la investigación disponible \cite{kitchenham2007guidelines}\cite{kitchenham2009slr} que sea relevante para una pregunta de investigación específica, un área temática o un fenómeno de interés. Se enfoca en el análisis detallado y profundo del tema en estudio. Su objetivo principal es sintetizar la evidencia disponible para tomar decisiones, a través de una metodología bien definida, imparcial, repetible y auditable. 
\section{Diferencias entre SMS y SLR}

Aunque ambos métodos son procesos sistemáticos y reproducibles, difieren principalmente en su alcance y profundidad, ya que el SMS se interesa en una comprensión de alto nivel como lo son el estado del arte y un estudio bibliográfico de un área más amplia, mientras que la SLR busca un estudio profundo de un tema en particular, para así responder preguntas puntuales y específicas.\\
Los resultados de una Revisión Sistemática de Literatura buscan sintetizar la evidencia de manera rigurosa para llegar a un análisis cualitativo y cuantitativo de los datos para responder a una hipótesis, en comparación con los resultados del Mapeo Sistemático que suelen ser gráficos de burbuja y mapas conceptuales.\\
Estas diferencias han sido descritas en detalle en la literatura metodológica sobre SLR y SMS en ingeniería de software. \cite{kitchenham2009slr}\\

\begin{table}[h!]
\centering
\caption{Principales diferencias entre la metodología SMS y la SLR}
\begin{tabular}{|p{2.1cm}|p{3.4cm}|p{3.4cm}|}
\hline
\textbf{Característica} & \textbf{Estudio de Mapeo Sistemático (SMS)} & \textbf{Revisión Sistemática de Literatura (SLR)} \\ \hline
\textbf{Objetivo} & Clasificar y categorizar un campo de investigación e identificar referencias disponibles. & Agregar y sintetizar evidencia para responder preguntas específicas. \\ \hline
\textbf{Alcance} & Amplio ya que busca mapear un área temática o fenómeno de interés de manera extensiva. & Específico: se centra de manera profunda en un tema o hipótesis muy puntual. \\ \hline
\textbf{Profundidad} & El análisis no es tan exhaustivo debido al volumen de estudios y el objetivo general. & Busca un entendimiento detallado y crítico de los hallazgos de los estudios primarios. \\ \hline
\textbf{Preguntas de Investigación} & Generales orientadas a tendencias, foros de publicación y distribución temática. & Puntuales y específicas como lo pueden ser factores de costo o riesgos. \\ \hline
\textbf{Evaluación de Calidad} &  Menos detallada ya que algunos criterios de calidad pueden ser menos aplicables que en una SLR. & Obligatoria y rigurosa, es esencial para asegurar que los hallazgos sean confiables y con menor sesgo. \\ \hline
\textbf{Resultados} & Mapas, tablas, gráficos de distribución y descubrimiento de brechas. & Una síntesis de evidencia que responde a la pregunta de investigación o hipótesis. \\ \hline
\textbf{Propósito estratégico} & Identificar áreas con exceso de publicaciones. & Proporcionar una base para la toma de decisiones informada o práctica basada en evidencia. \\ \hline
\end{tabular}
\end{table}
\section{Proceso metodológico}
El proceso para ambos se compone de tres fases: planificación, conducción y reporte, con diferencias fundamentales en el enfoque de cada fase para cada metodología.
\subsection{Planificación}
Durante la fase de planificación, se define la justificación y objetivos, se formulan las preguntas y se desarrolla un protocolo.\\ 

En el SMS, el protocolo\cite{definicion_sms} define:
\begin{itemize}
    \item Objetivo y preguntas de mapeo (amplias y exploratorias, enfocándose más en cantidad y tipo de estudios que en resultados detallados).
    \item Estrategias de búsqueda (bases de datos, palabras clave, etc.).
    \item Criterios de inclusión y exclusión de fuentes.
    \item Procedimiento de selección.
    \item Esquema de clasificación (define facetas: tema, tipo de contribución, tipo de investigación, contexto).
    \item Extracción de datos y análisis (formularios, plan de análisis de datos).
    \item Reporte (formato de presentación del mapa).\\
\end{itemize}

El protocolo del SLR\cite{kitchenham2009slr} se compone de:
\begin{itemize}
    \item Contexto, objetivo y preguntas de investigación o RQs (preguntas específicas buscando evidencias).
    \item Estrategias de búsqueda (bases de datos, palabras clave, etc.)
    \item Criterios de inclusión y exclusión de fuentes
    \item Procedimiento de selección
    \item Evaluación de calidad (checklists, uso de scores)
    \item Extracción y síntesis de datos (campos de extracción, tipo de síntesis, plan de análisis)
    \item Reporte (cronograma, versionamiento, plan de reporte)\\
\end{itemize}

\subsection{Conducción}
Este es el núcleo de ambas metodologías, aunque ambas tienen una estructura similar general, difieren significativamente en la profundidad y el tipo de datos que se procesan.\\
 
En el SMS, la conducción tiene como características que:
\begin{itemize}
    \item La investigación es mucho más amplia. El objetivo es cubrir toda la producción de un área temática o campo de estudio.
    \item La selección de estudios primarios puede ascender a los miles artículos los cuales se van filtrando a medida que avanza la investigación.
    \item La síntesis de datos es de alto nivel y se enfoca en la categorización y clasificación de los mismos.\\
\end{itemize}

Por otro lado, en la metodología SLR:
\begin{itemize}
    \item La investigación se enfoca en responder las preguntas de la investigación.
    \item La selección de estudios primarios es riguroso para filtrar estudios que no respondan directamente a la hipótesis o pregunta puntual.
    \item La síntesis de datos es profunda, así como cualitativa y cuantitativa.\\ 

\end{itemize}

\subsection{Reporte}
En la fase de reporte se describe el proceso y los resultados obtenidos pero la salida es diferente: una síntesis de evidencia en el SLR, y un mapa de terreno en el SMS.\\

En un SMS, el reporte típicamente utiliza las siguientes secciones para describir el paisaje:
\begin{itemize}
    \item Descripción del proceso.
    \item Presentación del esquema de clasificación.
    \item Resultados de mapeo (tablas y gráficos que muestren frecuencias y distribuciones).
    \item Identificación de tendencias y vacíos.\\
\end{itemize}

Por otro lado, el SLR está orientado a responder cada pregunta de investigación utilizando:
\begin{itemize}
    \item Descripción del proceso.
    \item Características de los estudios incluidos (tablas, scores).
    \item Síntesis de resultados.
    \item Discusión y conclusiones.\\
\end{itemize}
\section{Ejemplos}
En este enlace \cite{ejemplo_sms} se muestra un ejemplo de SMS realizado para el tema de bots en Ingeniería de Software, y en este otro \cite{ejemplo_slr} se presenta un ejemplo de SLR alrededor de calificaciones de aplicaciones dentro de las tiendas de aplicaciones (App Stores). \\
\section{Conclusiones}
 Tanto el Estudio de Mapeo Sistemático (SMS) como la Revisión Sistemática de Literatura (SLR) constituyen metodologías fundamentales y complementarias en el ámbito de la investigación científica. Aunque ambas comparten una estructura metodológica similar basada en tres fases principales (planificación, conducción y reporte), sus diferencias en términos de alcance, profundidad y objetivos las convierten en herramientas especializadas para propósitos distintos. El SMS se posiciona como una metodología ideal para obtener una visión panorámica de campos de investigación amplios, identificar tendencias emergentes y detectar brechas en el conocimiento, mientras que la SLR se enfoca en proporcionar respuestas específicas y detalladas a preguntas de investigación puntuales mediante un análisis riguroso y exhaustivo de la evidencia disponible.

% Bibliografía
\printbibliography

\end{document}